\section{Conclusion and Future Work}
\label{ch:conclusion}

Enhanced AI Racing achieved its central aim: it delivered a TORCS application in which contrasting autonomous drivers can be run, evaluated and explained through a research-informed novice-facing learning platform. The project implemented a meaningful progression from random and rule-based behaviour to map-aware control, Dyna-Q and N-step TD3. It also reduced deployment friction through a Windows package that contains TORCS, runtime policies and the application launcher.

The results show that the map-aware controller is a strong transparent baseline, completing 15 of 15 stored application attempts with a 92.056-second median completed lap. Agent~8 delivered the strongest observed final learned result, completing 36 of 40 repeated evaluations with a median completed-lap time of 83.038 seconds. These findings support the value of continuous deep reinforcement learning for this configuration, while the weaker Dyna-Q and scratch-TD3 outcomes show why such a conclusion cannot be assumed in advance.

The work remains limited to simulated one-track evaluation and functional rather than participant-based novice assessment. Future work should evaluate all selected agents in equal-sized frozen-policy batches; run policies on unseen tracks; perform an ablation that isolates racing-line input from reward and training differences; quantify variation with appropriate intervals; and conduct a small novice usability study. These steps would extend the current project from a strong evidence-aware simulator artefact towards a broader study of robust autonomous control and AI education.

The final lesson is that the strongest system outcome combines performance with inspectability. The project would be less useful if the fastest policy could be started only through private scripts or its failures were hidden in a checkpoint directory. Conversely, a polished interface without a reliable evidence trail would not meet the standard of a research artefact. Enhanced AI Racing connects these concerns: it preserves the technical detail necessary for scrutiny while giving a beginner a workable route into observing and discussing autonomous decision-making.

Within its scope, the work therefore makes three defensible claims. It demonstrates that several autonomous-driving approaches can be integrated behind one TORCS application and assessed with common outcome semantics. It provides evidence that the final sensor-only TD3 policy achieved the strongest observed learned performance on \texttt{g-track-3}, while the map-aware controller remained a stable and interpretable engineered benchmark. Finally, it demonstrates a packaged, evidence-aware and research-informed learning platform through which an AI newcomer can launch, observe and critically compare those findings without development setup. These are substantial outcomes, provided their limits are retained: they are not claims of real-world vehicle autonomy, universal algorithmic superiority or verified educational impact.

